\begin{eabstract}
	\addcontentsline{toc}{chapter}{ABSTRACT}
	The 3/2 stochastic volatility model, by virtue of its nonlinear diffusion structure, exhibits significant theoretical advantages in capturing extreme asset price volatility and fitting the volatility smile, making it a core tool for derivatives pricing in highly volatile markets. However, existing performance evaluations of pricing methods under this model are largely confined to idealized academic benchmark cases in stable markets, lacking empirical validation in actual extreme market environments. Consequently, the engineering applicability boundaries of these methods remain systematically unclarified. 

    To address this gap, this paper systematically constructs a comprehensive theoretical framework for nine European option pricing methods under the 3/2 model within a pure diffusion framework. These methods are categorized into three major paradigms: time-stepping conditional Monte Carlo methods, exact and near-exact terminal simulation methods, and Fourier frequency-domain pricing methods. The implementation logic, error characteristics, and convergence properties of each paradigm are rigorously derived. Building upon this, through seven sets of numerical experiments covering both conventional markets and extreme parameter environments, this study quantitatively compares the pricing accuracy, computational efficiency, and full-scenario numerical stability of these methods, thereby establishing a clear hierarchical performance framework.

    Furthermore, taking the extreme Bitcoin flash crash of October 2025 as a case study, this paper develops a complete empirical framework consisting of ``long-term fixed-parameter time-series estimation — time-varying parameter cross-sectional calibration — pricing method re-evaluation in extreme scenarios.'' This validates the 3/2 model's adaptability to extreme cryptocurrency market dynamics, identifies the event-driven time-varying patterns of core model parameters throughout the crash cycle, and executes a performance re-evaluation of the pricing methods under genuine extreme market conditions. 

    The findings indicate that the Fourier-Cosine (Fourier-Cos) method is the globally optimal scheme for European option pricing, achieving exceptional pricing precision, numerical stability, and computational efficiency across all scenarios. The Poisson-Gamma and Quasi-Exponential (QE) near-exact stepping methods are the sole robust solutions across all scenarios within the Monte Carlo framework, offering a reliable implementation pathway for pricing path-dependent derivatives and dynamic hedging. Conversely, the nominally ``exact'' terminal simulation methods can only be theoretically validated under idealized benchmark parameters; they prevalently experience numerical overflows and systematic pricing deviations in extreme market environments, thus lacking suitability for practical engineering applications. 

    This research bridges the existing disconnect between theoretical simulations and real-world markets in 3/2 model pricing literature, establishes method selection guidelines for diverse application scenarios, and provides systematic theoretical support and empirical foundations for derivatives pricing, model calibration, and risk management in highly volatile markets.
\end{eabstract}

% vim:ts=4:sw=4
